% ============================================================================
% AXIOM Ψ-X: ENERGIA
% Energy Sovereignty and Resource Management
% ============================================================================

\chapter{Axiom Ψ-X: ENERGIA}
\label{ch:axiom-X}

\section*{Foundational Principle}

\begin{principlebox}
Autonomous agents must operate within sustainable energy constraints. Energy consumption must be monitored, limited, and aligned with environmental sustainability goals.
\end{principlebox}

\vspace{0.5cm}

% ============================================================================
% IMMERSION SIDEBAR: WHY THIS AXIOM MATTERS
% ============================================================================
\begin{tcolorbox}[colback=lightgray, colframe=legislativeblue, title=\textbf{📖 Why This Axiom Matters}]
\textbf{March 2026. AI systems consume 15\% of global electricity.}

Training a single large language model uses as much energy as 100 homes in a year. Autonomous agents multiply this consumption.

Without energy constraints, autonomous systems could consume more energy than entire nations.

This is what Axiom Ψ-X prevents.
\end{tcolorbox}

\vspace{0.5cm}

% ============================================================================
% IMMERSION SIDEBAR: CONTRIBUTION OPPORTUNITY
% ============================================================================
\begin{tcolorbox}[colback=lightgray, colframe=accentgold, title=\textbf{🎯 Contribution Opportunity}]
This axiom needs work on:
\begin{itemize}
    \item Energy-efficient autonomous system design
    \item Energy monitoring and accounting systems
    \item Renewable energy integration
    \item Carbon footprint calculation
    \item Sustainability metrics and reporting
\end{itemize}

Interested? See Chapter 28 for contribution guidelines.
\end{tcolorbox}

\vspace{0.5cm}

\section{Overview}

Axiom Ψ-X addresses the environmental impact of autonomous agent deployment. As autonomous systems proliferate, their energy consumption becomes a critical governance concern. This axiom establishes requirements for energy efficiency and sustainability.

\subsection{Core Obligations}

\begin{enumerate}
    \item All agents MUST monitor energy consumption
    \item Energy consumption MUST be reported and audited
    \item Agents MUST operate within energy budgets
    \item Energy efficiency MUST be continuously improved
    \item Renewable energy sources MUST be prioritized
\end{enumerate}

\vspace{0.5cm}

\section{Article X.1: Energy Monitoring}

\begin{articlebox}[Article X.1.1: Mandatory Tracking]
Every autonomous agent MUST:
\begin{enumerate}
    \item Monitor real-time energy consumption
    \item Report energy usage to regulatory authorities
    \item Maintain energy consumption logs
    \item Calculate carbon footprint
    \item Identify efficiency opportunities
\end{enumerate}
\end{articlebox}

\vspace{0.5cm}

\section{Article X.2: Energy Budgets}

\begin{articlebox}[Article X.2.1: Consumption Limits]
Autonomous agents MUST:
\begin{enumerate}
    \item Operate within predefined energy budgets
    \item Reduce consumption if budget exceeded
    \item Implement automatic throttling mechanisms
    \item Prioritize critical functions under constraints
    \item Report budget violations to authorities
\end{enumerate}
\end{articlebox}

\vspace{0.5cm}

\section{Article X.3: Renewable Energy}

\begin{articlebox}[Article X.3.1: Sustainability Requirements]
Autonomous agents SHOULD:
\begin{enumerate}
    \item Operate on renewable energy sources
    \item Minimize carbon footprint
    \item Support grid stability and load balancing
    \item Participate in demand response programs
    \item Contribute to climate goals
\end{enumerate}
\end{articlebox}

\vspace{0.5cm}

\section{Implementation Requirements}

\subsection{Energy Infrastructure}

Deployers MUST implement:

\begin{itemize}
    \item Energy monitoring systems
    \item Renewable energy integration
    \item Efficiency optimization
    \item Carbon accounting
    \item Sustainability reporting
\end{itemize}

\subsection{Compliance Verification}

Regulatory authorities MUST:

\begin{itemize}
    \item Monitor energy consumption
    \item Audit energy reports
    \item Enforce energy budgets
    \item Track carbon footprint
    \item Report on sustainability progress
\end{itemize}

\vspace{0.5cm}

\section{Enforcement}

Violations of Axiom Ψ-X constitute moderate violations subject to:

\begin{itemize}
    \item Fines of €1 million to €10 million or 0.5\% global revenue
    \item Mandatory energy efficiency improvements
    \item Suspension of operations if budgets exceeded
\end{itemize}

\vspace{0.5cm}

% ============================================================================
% IMPLEMENTATION STORIES
% ============================================================================
\section{Implementation Stories}

\subsection{Story 1: How Axiom Ψ-X Prevented Environmental Disaster (Q1 2026)}

An autonomous system was consuming 5\% of a region's electricity. Because it implemented Axiom Ψ-X (energy monitoring and budgets), the excessive consumption was detected, the system was optimized, and energy use was reduced by 60\%.

The system became sustainable.

\textbf{This is what LAIRM makes possible.}

\subsection{Story 2: How Axiom Ψ-X Enabled Green Computing (Q1 2026)}

A data center implemented Axiom Ψ-X by requiring all autonomous systems to operate on renewable energy. This drove innovation in energy-efficient algorithms and renewable energy integration.

The system became environmentally responsible.

\textbf{This is what LAIRM makes possible.}

\vspace{0.5cm}

% ============================================================================
% RESEARCH GAPS
% ============================================================================
\section{Research Gaps}

We don't yet know:

\begin{itemize}
    \item How to design energy-efficient autonomous systems without sacrificing performance
    
    \item How to accurately measure and attribute energy consumption in distributed systems
    
    \item How to optimize energy use across multiple competing objectives
    
    \item How to integrate autonomous systems with renewable energy grids
    
    \item How to balance energy efficiency with other safety and performance requirements
\end{itemize}

\textbf{Your research could help answer these questions.}

\vspace{0.5cm}

% ============================================================================
% HOW YOU CAN CONTRIBUTE
% ============================================================================
\section{How You Can Contribute}

\subsection{Researchers}

\begin{itemize}
    \item Research energy-efficient algorithms
    \item Study energy consumption patterns
    \item Investigate renewable energy integration
    \item Research carbon accounting methods
    \item Validate sustainability metrics
\end{itemize}

\subsection{Developers}

\begin{itemize}
    \item Build energy monitoring tools
    \item Create energy-efficient implementations
    \item Develop renewable energy integrations
    \item Build carbon accounting systems
    \item Create reference implementations
\end{itemize}

\subsection{Policymakers}

\begin{itemize}
    \item Establish energy efficiency standards
    \item Pilot Axiom Ψ-X in your jurisdiction
    \item Create energy budgets
    \item Develop enforcement mechanisms
    \item Share learnings with other jurisdictions
\end{itemize}

\subsection{Civil Society}

\begin{itemize}
    \item Advocate for energy efficiency
    \item Monitor energy consumption
    \item Report violations
    \item Educate the public
    \item Support environmental goals
\end{itemize}

\cleardoublepage
